% Practical explanation added to Appendix C; no fitted results.
\renewcommand{\backupcode}{C5}
\begin{frame}[more-space,label=backup-C5]{C5\enspace|\enspace Physics in practice: fit the measurements}
  \small
  \dataheading{1. Choose the polynomial coefficients}
  \begin{beamercolorbox}[sep=2.5mm,wd=\textwidth]{data panel}
    \[
      \widehat T(\mathbf x;\mathbf c)
      =\sum_{|\alpha|\leq p}c_\alpha\Psi_\alpha(\boldsymbol\xi),
      \qquad \mathbf x=(P,T_{\mathrm{inlet}},F).
    \]
    The polynomial terms $\Psi_\alpha$ are fixed. The unknowns are the coefficients $\mathbf c$.
    \par\smallskip
    Inputs are scaled to $\boldsymbol\xi\in[-1,1]^3$ using training ranges only.
  \end{beamercolorbox}
  \par\vfill
  \dataheading{Ordinary PCE: minimize squared prediction errors}
  \[
    \min_{\mathbf c}\;\sum_{n\in\mathrm{training}}
      \left[\widehat T(\mathbf x_n;\mathbf c)-y_n\right]^2.
  \]
  $y_n$: largest of the 18 CPU temperature averages for measured interval $n$.
  \par\vfill
  \begin{beamercolorbox}[sep=2mm,wd=\textwidth]{data warning}
    \textbf{PI-PCE uses the same measurements and objective,}
    while restricting the polynomial to expected heating and cooling directions.
  \end{beamercolorbox}
\end{frame}

\renewcommand{\backupcode}{C6}
\begin{frame}[more-space,label=backup-C6]{C6\enspace|\enspace Physics in practice: constrain the slopes}
  \small
  \dataheading{2. Add monotonicity constraints}
  \begin{beamercolorbox}[sep=2mm,wd=\textwidth]{data panel}
    \[
      \frac{\partial\widehat T}{\partial P}\geq0,
      \qquad \frac{\partial\widehat T}{\partial T_{\mathrm{inlet}}}\geq0,
      \qquad \frac{\partial\widehat T}{\partial F}\leq0.
    \]
    Change one input while holding the other two fixed:
    more power / warmer inlet cannot cool; more fan cooling cannot heat.
  \end{beamercolorbox}
  \par\vfill
  \dataheading{A simple linear example}
  Written in scaled coordinates, with slopes $a_P,a_T,a_F$:
  \[
    \widehat T=a_0+a_P\xi_P+a_T\xi_T+a_F\xi_F
    \quad\Longrightarrow\quad
    a_P\geq0,\quad a_T\geq0,\quad a_F\leq0.
  \]
  {\footnotesize\color{black!60}For normalized linear Legendre terms, each slope is
    $\sqrt3$ times its corresponding PCE coefficient; the sign is unchanged.\par}
\end{frame}

\renewcommand{\backupcode}{C7}
\begin{frame}[more-space,label=backup-C7]{C7\enspace|\enspace Physics in practice: check eight corners}
  \small
  \dataheading{3. With quadratic terms, slopes depend on the operating point}
  For example, using ordinary polynomial notation:
  \[
    \widehat T=\cdots+a_F\xi_F+a_{PF}\xi_P\xi_F+a_{FF}\xi_F^2
    \quad\Longrightarrow\quad
    \frac{\partial\widehat T}{\partial\xi_F}
      =a_F+a_{PF}\xi_P+2a_{FF}\xi_F.
  \]
  \par\vfill
  \begin{beamercolorbox}[sep=2.5mm,wd=\textwidth]{data panel}
    \dataheading{Eight corners $\times$ three derivative signs = 24 inequalities}
    \begin{itemize}
      \item Check all corners of the training-input box $[-1,1]^3$.
      \item For total degree 1 or 2, every derivative is affine:
        its minimum / maximum occurs at a corner.
      \item The signs therefore hold throughout the box, to solver tolerance.
    \end{itemize}
  \end{beamercolorbox}
\end{frame}

\renewcommand{\backupcode}{C8}
\begin{frame}[more-space,label=backup-C8]{C8\enspace|\enspace Physics in practice: solve the constrained fit}
  \small
  \dataheading{4. Adjust all coefficients together}
  \begin{beamercolorbox}[sep=2mm,wd=\textwidth]{data panel}
    \[
      \boxed{\min_{\mathbf c}\|A\mathbf c-\mathbf y\|_2^2
        \quad\text{subject to}\quad B\mathbf c\geq0}
    \]
    $A$: polynomial terms evaluated at measured inputs.
    \par\smallskip
    $B$: signed derivatives of those terms evaluated at the eight corners.
  \end{beamercolorbox}
  \par\vfill
  \begin{enumerate}
    \item Compute the ordinary least-squares fit.
    \item If it satisfies all signs, retain the same coefficients.
    \item Otherwise, use SciPy SLSQP to minimize error subject to the constraints;
      check convergence and final derivative signs.
  \end{enumerate}
  \par\vfill
  \begin{beamercolorbox}[sep=2mm,wd=\textwidth]{data warning}
    \textbf{Physical guidance where data are sparse.} The signs remain assumptions
    to validate, especially for the raw fan signal. They do not establish prediction
    accuracy, safety or behavior outside the training box.
  \end{beamercolorbox}
  \datasource{Solver: \href{https://docs.scipy.org/doc/scipy/reference/optimize.minimize-slsqp.html}{SciPy SLSQP documentation}.
    Basis and earlier details: \hyperlink{backup-C1}{\textbf{C1}}, \hyperlink{backup-C3}{\textbf{C3}}.}
\end{frame}
